%% Ethics Statement for the CBM submission. Build: latexmk -pdf ethics_statement.tex
\documentclass[11pt]{article}
\usepackage[T1]{fontenc}
\usepackage{lmodern}
\usepackage[margin=2.5cm]{geometry}
\usepackage[hidelinks]{hyperref}
\usepackage{xurl}
\usepackage{orcidlink}
%% Target journal and author identifiers shared by every submission document.
\providecommand{\targetjournal}{Biomedical Signal Processing and Control}
\providecommand{\authororcid}{0009-0003-7760-0522}

\pagestyle{empty}
\setlength{\parindent}{0pt}
\setlength{\parskip}{8pt}
\begin{document}

{\Large\bfseries Ethics Statement}

\textbf{Manuscript title:} Leakage-controlled evaluation of stacking and recursive feature elimination for diabetes outcome prediction on the PIMA dataset

\textbf{Journal:} \targetjournal

%% TODO(author): confirm name and affiliation; check whether ANITS policy requires a formal exemption letter.
\textbf{Corresponding author:} Subrahmanyam Tanala, Department of Electrical and Electronics Engineering, Anil Neerukonda Institute of Technology and Sciences (ANITS), Visakhapatnam, Andhra Pradesh, India (\href{mailto:subrahmanyam.eee@anits.edu.in}{subrahmanyam.eee@anits.edu.in}); ORCID \href{https://orcid.org/\authororcid}{\authororcid}

\bigskip
\textbf{Human participants and data.} This study is a secondary analysis of the publicly available PIMA Indians Diabetes Database, originally collected by the National Institute of Diabetes and Digestive and Kidney Diseases and distributed through the UCI Machine Learning Repository. The dataset contains 768 records with eight clinical measurements and a binary outcome. It contains no names, dates, locations or other direct identifiers. The replication analysis used the public-use files of the US National Health and Nutrition Examination Survey (NHANES) 2013--2014, released by the National Center for Health Statistics (NCHS) without direct identifiers. NHANES protocols were approved by the NCHS Research Ethics Review Board, and all participants gave written informed consent at the time of data collection. The author had no access to identifiable information and no contact with participants.

\textbf{Ethical approval.} No new data were collected from human participants, no intervention was performed, and only de-identified public data were analysed. Ethical approval was therefore not required for this study.

\textbf{Informed consent.} Not applicable to this secondary analysis. The study used de-identified, publicly available secondary data, and no individual can be identified from the data or from the results reported.

\textbf{Animal research.} Not applicable. No animals were involved.

\textbf{Data use.} Both datasets were used solely for methodological research, consistent with their public distribution for research purposes and, for NHANES, with the NCHS data use restrictions (no attempt to identify participants). The exact files used, their sources and their checksums are documented in the supplementary material, and all code is publicly available.

\textbf{Responsible reporting.} Because the dataset represents a single historical cohort of women of Pima heritage, the manuscript explicitly states that the results are an internal methodological benchmark and must not be used for clinical decisions without external validation.

\vspace{2cm}
\begin{tabular}{@{}p{7cm}p{6cm}@{}}
\rule{6cm}{0.4pt} & \rule{5cm}{0.4pt} \\
Subrahmanyam Tanala & Date
\end{tabular}

\end{document}
